\section{Limitations and Claims This Paper Does Not Make}
\label{sec:limitations}

\paragraph{Limitations in scope.}
We validate our manipulation checks on another pretraining family (\S\ref{sec:replication}).
However, our benchmark comparisons, leave-one-out analysis, and single call comparison are only
performed within a single pretraining family with a similar level of capability. Validation on a
second family only establishes replication, not population-level generalization. We also only
examine two benchmarks, both in English.

\paragraph{Limitations in measurement.}
All our arms return the same format; this means our manipulation checks are measuring the relative
difference rather than the existence of these factors, resulting in lower contrast relative to an
experiment where the null condition doesn't return anything. The concentration and sourced claim
manipulation checks were designed as two-factorial but effectively function as a one-factorial
design since both conditions are saturated; this collapses their corresponding statistics. We only
filter refusals in TruthfulQA, where there is no individual question contributing to the accuracy
comparisons (Appendix~\ref{app:refusal}). While the role priming could influence whether the model
refuses to answer at all, it did not affect that measured benchmark in our experiment.

\paragraph{Limitations in provenance.}
We directly state, but cannot verify, two pieces of information (\S\ref{sec:models}). First, we
state the model variant but it cannot be deduced from the serving name. Second, we state the serving
path but it cannot be deduced from the manifest since our scraper does not record provider. We
conducted an experiment on a 47-question test dataset before our registered analysis; no out-dated
correction affected a registered quantity, and we ask the reader to consider any analysis they
believe affected to be exploratory. Our analysis of failure modes in attribution are only evaluated
on simulated data in Appendix~\ref{app:attribution-failures} to serve as evidence of existence and
possible mechanisms rather than rates of failure in public models.

\paragraph{What this does not claim: beating baselines.}
This scoped experiment isn't powered for that. We also don't claim to have found correspondences
between LLM capabilities and clinical constructs; this was only inspired by that analogy. We do not
claim to be the first to do this work; we just note that role checks are rarely done. We don't track
how often attribution errors occur. We did not test H1 (not powered by design) nor H2 (lost its
target when we applied the two benchmark simplification and did not choose an alternative contrast
such as HaluEval); provide cost-matched comparisons of council vs single agent: we recorded a
same-model single-call arm, 6$\times$ cheaper; to calibrate: we do not propose a calibration
procedure, Appendix~\ref{app:confidence} is a gap analysis rather than a solution. Finally, we do
not claim that capability is broken when the constraint that we provided is lifted; this requires
training tests we didn't run.

\paragraph{Solely for research.}
Nothing here licenses deployment. The council described in this paper must not be used for medical,
legal, financial or other consequential decisions on the basis of benchmark accuracy.
